\chapter*{Conclusion générale}
\addcontentsline{toc}{chapter}{Conclusion générale}

\section*{Conclusion}

La réalisation de ce projet de fin d'année, intitulé \textbf{\nomapplication}, nous a permis de mener un développement web complet, de la conception UML jusqu'à une application fonctionnelle et testée. Destinée à la gestion d'une agence bancaire, cette plateforme répond aux exigences du cahier des charges et intègre des extensions métier : paiement de factures, commande de chéquier, notifications et portail client.

\section*{Bilan}

\begin{itemize}[leftmargin=*]
    \item \textbf{Conception} : diagrammes UML complets (cas d'utilisation, classes, séquences, activité, MCD/MLD).
    \item \textbf{Implémentation} : 11 migrations Flyway, 9 entités JPA, architecture en couches Spring Boot.
    \item \textbf{Fonctionnalités} : gestion clients/comptes, opérations financières contrôlées, admin, audit, reporting, reçus PDF.
    \item \textbf{Extensions} : factures, chéquier, notifications bidirectionnelles, portail client lecture seule.
    \item \textbf{Qualité} : 80 tests d'intégration, seed de démo modulaire, documentation technique complète.
    \item \textbf{Interface} : charte AmanaBank cohérente (landing, agence, portail).
\end{itemize}

Les objectifs fixés en introduction ont été atteints : centralisation des données, respect des règles métier, traçabilité et expérience utilisateur adaptée à chaque rôle.

\section*{Difficultés rencontrées}

\begin{itemize}[leftmargin=*]
    \item \textbf{Configuration PostgreSQL} : conflit de ports locaux (5432 occupé) — résolu en utilisant le port 5433.
    \item \textbf{Cohérence transactionnelle} : gestion des soldes avec verrouillage optimiste (\texttt{@Version}) et \texttt{@Transactional}.
    \item \textbf{Authentification dual} : un seul formulaire de login pour personnel et clients — résolu via \texttt{UserDetailsServiceImpl} composite.
    \item \textbf{Portail sur base existante} : synchronisation des accès client via \texttt{DemoPortalSync} au démarrage.
    \item \textbf{Workflow chéquier} : distinguer service administratif (sans mouvement financier) des transactions.
\end{itemize}

\section*{Perspectives d'évolution}

\begin{itemize}[leftmargin=*]
    \item Export PDF des \textbf{relevés de compte} (reçus déjà en PDF).
    \item API REST documentée (Swagger) pour intégrations futures.
    \item Pipeline CI/CD (GitHub Actions) et déploiement conteneurisé.
    \item Enrichissement du portail client (demandes en ligne, messagerie).
    \item Authentification à deux facteurs (2FA) pour le personnel.
    \item Tableau de bord analytique avancé (graphiques, exports Excel).
\end{itemize}

\vspace{0.5cm}
En conclusion, \nomapplication\ constitue une solution concrète et défendable pour la digitalisation d'une agence bancaire. Ce PFA nous a permis de consolider nos compétences en conception orientée objet, développement Java/Spring et travail en binôme, et constitue une base solide pour nos futurs projets professionnels.

\vspace{1cm}
\begin{flushright}
    \etudianteun\\
    \etudiantedeux
\end{flushright}
